% ===================================================================
%  तक्ता क्रमांक ९ — डोंगर. रान.
%
%  Traduction, faite en 2026, en marathi d'aujourd'hui, sur l'ido de
%  texto/io. Regles de la colonne : voir 10-takta-01.tex. Deux
%  scenes, deux numerotations : les cles portent c1 et c2.
%
%  LE MARATHI EST CHEZ LUI DANS CE TABLEAU, ET C'EST GEOGRAPHIQUE.
%  Le Maharashtra est un pays de montagne : les Ghats, les forts, les
%  cols. La ou l'ido compose, le marathi a le mot tout fait, et c'est
%  le mot qu'on lit tous les jours sur une carte :
%
%    किल्ला (8)   la forteresse    घाट (27)    le defile
%    दरी (25)     la vallee        घळ (32)     le ravin
%    कडा          le precipice     माथा (2)    la crete
%    धबधबा (10)   la cascade       पायवाट (53) le sentier
%    गुहा (56)    la grotte        खडक (45)    les rochers
%    वाटाड्या (48) le guide        डोंगरवासी (38) le montagnard
%    सडक (22)     la chaussee      बांध (23)   le talus
%
%  ET LA LIEUE TOMBE EXACTEMENT. L'ido n'a pas de mot pour elle et
%  ecrit « quaropa kilometri », des kilometres quadruples. Le marathi
%  dit कोस — l'unite de distance marathe, qui vaut a peu pres ces
%  quatre kilometres-la. La ou l'ido devait calculer, une langue
%  avait deja compte.
%
%  DEUX SERPENTS, ET LE MARATHI FAIT LA MEME OPPOSITION QUE L'IDO.
%  Le texte dit : je l'ai d'abord cru vipero, mais ce n'etait qu'un
%  simple kolubro. Le marathi ecrit फुरसे — la vipere a ecailles
%  scieuses, qui tue — contre धामण, la couleuvre inoffensive que tout
%  le monde connait. Deux mots courants, la meme peur et le meme
%  soulagement.
%
%  LES ARBRES ET LES BETES, MOT POUR MOT : देवदार le cedre, चिनार le
%  peuplier, भूर्ज le bouleau, राळ la resine, शेवाळ la mousse,
%  सुतारपक्षी le pic, राजहंस le cygne, सांबर le cerf, भेकर le
%  chevreuil, रानडुक्कर le sanglier, खार l'ecureuil, वारूळ la
%  fourmiliere, गोगलगाय l'escargot, अळंबी le champignon, नेचा la
%  fougere, सदाफुली la pervenche.
%
%  ET LES EMPRUNTS RESTENT DES EMPRUNTS, comme au tableau 8 : la
%  marmotte, le chamois, la bruyere, l'if, le hetre, l'orme, la pie
%  ne sont pas d'ici. On les nomme par l'emprunt.
%
%  UNE INVERSION PREVENUE — le troupeau qui pait dans le paturage —
%  ET SEPT QUI ONT PASSE, toutes vues par renvoji.py.
%
%  MAIS LE PLUS UTILE EST VENU DE kolonoj.py, ET C'EST LA PREMIERE
%  FOIS. Son controle de bloc ampute a signale que t09-c1-05-1 ne
%  faisait que 37 % de la longueur mediane de la colonne. La raison :
%  le second alinea du cinquieme n'a PAS de %%K dans le fac-simile —
%  il appartient a c1-05-1, que le changement de feuillet a coupe —
%  et j'en avais fait un bloc a lui, sous une cle qui n'existe pas.
%  renvoji.py disait seulement « cle inconnue » du bloc en trop ; il
%  ne pouvait pas dire que l'AUTRE bloc, celui qui existe, avait
%  perdu la moitie de son texte. Le controle ne de la chirurgie
%  marathe du tableau 4 vient de payer sa place.
% ===================================================================

%%K t09-apar-1 apar
\VUsaut{4.0mm}
\VUtitre{15.0pt}{60}{तक्ता क्रमांक ९}
\VUsaut{3.0mm}

%%K t09-tit-1 sub
\VUcentre{12.0pt}{0}{\textit{पहिला प्रवेश.}}
\VUsaut{3.0mm}

%%K t09-tit-2 sub
\VUpk{11.4pt}{40}{डोंगर. --- स्नानाचे शहर.}
\VUsaut{3.0mm}

%%K t09-c1-01-1 p
१. --- आठ दिवस झाले मी प्रात्झमध्ये आहे. पिरेनीजच्या त्या अद्भुत \VUgras{पर्वतरांगेसमोर}
\textsuperscript{(1)} उभा राहिलो तेव्हा मला जे कौतुक वाटले ते शब्दांत सांगता येणार नाही — जिचा
\VUgras{माथा} \textsuperscript{(2)} निळ्या आकाशावर एखाद्या अवाढव्य तोरणासारखा दिसतो.

%%K t09-c1-02-1 p
२. --- पिरेनीजची \VUgras{शिखरे} \textsuperscript{(3)} एवढी \VUgras{उंची} गाठतात याची मी
कल्पनाही केली नव्हती, आणि त्या देखण्या \VUgras{हिमनद्यांसमोर} \textsuperscript{(5)} आणि बर्फाळ
\VUgras{डोंगरटोकांसमोर} \textsuperscript{(4)} मी थक्क होऊन उभा राहिलो — ज्यांच्यावरून एवढे
भयंकर \VUgras{हिमप्रपात} घसरत येतात.

%%K t09-tit-3 sub
\VUsaut{3.0mm}
\VUpk{10.2pt}{40}{डोंगरी निसर्गदृश्य.}
\VUsaut{3.0mm}

%%K t09-c1-03-1 p
३. --- मी पोहोचल्याच्या दुसऱ्याच दिवशी मी प्रात्झजवळच्या त्या जुन्या विझलेल्या
\VUgras{ज्वालामुखीकडे} \textsuperscript{(6)} गेलो. \VUgras{विवराच्या} रुक्ष आणि उघड्या कडांवर,
कित्येक शतकांपूर्वी \VUgras{डोंगराच्या उतारावरून} \textsuperscript{(7)} वाहिलेल्या
\VUgras{लाव्ह्याच्या} खुणा अजूनही स्पष्ट दिसतात. एका \VUgras{उतारावर} मला त्या लाव्ह्याचे काही
तुकडे सापडले.

%%K t09-c1-04-1 p
४. --- प्रात्झ सरहद्दीवर वसले आहे, आणि तो \VUgras{किल्ला} \textsuperscript{(8)} — जो फ्रान्सला
स्पेनपासून वेगळा करणाऱ्या \VUgras{घाटाचे} \textsuperscript{(27)} रक्षण करतो — ऱ्हाइनकाठच्या
त्या जुन्या किल्ल्यांची अगदी आठवण करून देतो, जे आपल्या खडकाच्या माथ्यावरून \VUgras{भक्ष्याची}
टेहळणी करत असल्यासारखे दिसतात.

%%K t09-c1-05-1 p
५. --- एक अवाढव्य \VUgras{गरुड} \textsuperscript{(9)}, ज्याचे घरटे \VUgras{किल्ल्यापासून}
\textsuperscript{(8)} फार दूर नाही, तो माझे लक्ष वारंवार वेधून घेतो. तो \VUgras{शिकारी पक्षी},
ज्याची \VUgras{चोच} बळकट आहे, आपल्या पिल्लांसाठी अनेकदा एखादे कोकरू किंवा करडू आणतो, जे तो
आपल्या \VUgras{नखांत} धरतो. त्याच्या \VUgras{पंखांचा} विस्तार एवढा आहे की तो आपले जड ओझे घेऊन
बराच वेळ तरंगत उडू शकतो.


%%K t09-tit-4 sub
\VUsaut{3.0mm}
\VUpk{10.2pt}{40}{भटकंती करणारा.}
\VUsaut{3.0mm}

%%K t09-c1-06-1 p
६. --- तिथला सर्वात आल्हाददायक फेरफटका म्हणजे « काऊ »च्या \VUgras{धबधब्यापर्यंतची}
\textsuperscript{(10)} भटकंती. \VUgras{जलप्रपात} भोवऱ्यात कोसळतो आणि पुढे शहराच्या पश्चिमेला
असलेल्या \VUgras{फरच्या बनात} \textsuperscript{(11)} एक सुंदर \VUgras{ओढा} होऊन नाहीसा होतो.
या बनातून आम्ही \VUgras{हवामान वेधशाळेपर्यंत} \textsuperscript{(12)} चढून गेलो, आणि
\VUgras{सज्ज्याच्या} माथ्यावरून सगळ्या प्रदेशाचा \VUgras{दृश्यपट} आम्हाला दिसला.
\VUgras{निसर्गदृश्य} अद्भुत आहे, आणि \VUgras{निसर्गाने} आपल्याजवळचे सगळे खजिने या प्रदेशावर
उधळले आहेत असे वाटते.

%%K t09-c1-07-1 p
७. --- उतरण्यासाठी आम्ही \VUgras{रज्जुमार्ग} \textsuperscript{(13)} वापरला आणि एका लहानशा
\VUgras{स्थानकावर} थांबलो — प्रात्झच्या सरदारांनी पूर्वी वस्ती केलेल्या जुन्या
\VUgras{सरंजामी किल्ल्याच्या} \VUgras{अवशेषांजवळ} \textsuperscript{(14)}.

%%K t09-c1-08-1 p
८. --- बिचारा किल्ला! खाली ते नखरेल \VUgras{स्नानाचे शहर} \textsuperscript{(15)} पाहून तो काय
बरे विचार करत असेल — जे आता फॅशनेबल \VUgras{उष्णजल केंद्र} झाले आहे?

%%K t09-c1-08-2 p
\VUgras{हवामान} इतके आल्हाददायक आहे, \VUgras{खनिज पाणी} इतके गुणकारी आहे, की \VUgras{आरोग्यधामात} \textsuperscript{(17)} घेतले जाणारे \VUgras{उपचार} म्हणजे खरा आनंद असतो.

%%K t09-tit-5 sub
\VUsaut{3.0mm}
\VUpk{10.2pt}{40}{आल्हाददायक स्नानशहर.}
\VUsaut{3.0mm}

%%K t09-c1-09-1 p
९. --- शिवाय, मुक्काम सुखाचा व्हावा म्हणून सगळे काही केले आहे. रेल्वे शक्य व्हावी म्हणून मोठा
\VUgras{रेल्वेपूल} \textsuperscript{(16)} बांधला; \VUgras{उष्णजल संस्थेची} \VUgras{बाग}
\textsuperscript{(18)}, जिथे \VUgras{मैफली} होतात, ती मोठ्या मोहक रीतीने मांडली आहे. कित्येक
रेखीव « \VUgras{बंगले} » \textsuperscript{(19)} बांधले आहेत — \VUgras{कॅसिनोभोवती}
\textsuperscript{(21)} —, आणि फार चांगली राखलेली \VUgras{सडक} \textsuperscript{(22)} ये-जा
खूपच सोपी करते.

%%K t09-c1-10-1 p
१०. --- एक साधा \VUgras{बांध} \textsuperscript{(23)} \VUgras{रस्त्याला} \textsuperscript{(20)}
खऱ्या \VUgras{दरीपासून} \textsuperscript{(25)} वेगळा करतो, जिच्या तळाशी एक रेखीव लहानसे
\VUgras{तळे} \textsuperscript{(26)} पडले आहे, जे एका नितळ \VUgras{ओढ्याला}
\textsuperscript{(24)} पाणी पुरवते.

%%K t09-c1-11-1 p
११. --- दरीच्या पलीकडच्या बाजूला एक \VUgras{लोहखाण} \textsuperscript{(28)} चालवली जाते.
\VUgras{खाणकामगार} \VUgras{खाणविहिरीत} उतरताना मी अनेकदा पाहायला गेलो, आणि ज्या
\VUgras{बोगद्यात} ते आपला दिवस काढतात, \VUgras{धातू} असलेला \VUgras{कच्चा खनिज} काढत, त्या
बोगद्यातही मी शिरलो.

%%K t09-c1-12-1 p
१२. --- खाणीतून निघणाऱ्या संपत्तीचा लगेच उपयोग व्हावा म्हणून खाणीजवळ एक मोठा
\VUgras{ओतकारखाना} बांधला आहे. इथे मुबलक असलेल्या \VUgras{दगडी कोळशाने} तापवलेली
\VUgras{मोठी भट्टी} \textsuperscript{(29)} झटपट आणि स्वस्तात \VUgras{बीड} मिळवणे शक्य करते.

%%K t09-c1-13-1 p
१३. --- खाणीपासून फार दूर नाही तिथे एक \VUgras{दगडखाण} \textsuperscript{(30)} खणली जाते.
म्हातारा \VUgras{दगडफोड्या} आपल्या \VUgras{टिकावाने} ना \VUgras{ग्रॅनाइट} कापतो, ना
\VUgras{पॉर्फिरी}, ना \VUgras{वालुकाश्मसुद्धा}, तर घरे बांधण्यासाठी साधे \VUgras{चिरे}.

%%K t09-c1-14-1 p
१४. --- त्या प्रदेशाच्या सर्वात सुंदर शोभांपैकी एक म्हणजे \VUgras{फरची झाडे}
\textsuperscript{(31)}. सगळीकडे — \VUgras{घळीच्या} \textsuperscript{(32)} तळाशी,
\VUgras{खळाळत्या ओहोळाच्या} \textsuperscript{(33)}, \VUgras{खोल दरीच्या}
\textsuperscript{(34)} किंवा \VUgras{कड्याच्या} काठावर — त्यांच्या बारीक \VUgras{सुया} हिरव्या
छटा उधळतात.

%%K t09-tit-6 sub
\VUsaut{3.0mm}
\VUpk{10.2pt}{40}{पायी भटकंती.}
\VUsaut{3.0mm}

%%K t09-c1-15-1 p
१५. --- सुट्टीचा फायदा घेऊन मी \VUgras{लांबवर} फिरतो. डोंगरावर वळणे घेत जाणारे \VUgras{रस्ते}
\textsuperscript{(35)} अंतराची पर्वा न करता कित्येक \VUgras{किलोमीटर} आणि अनेकदा कित्येक
\VUgras{कोस} तुडवणे शक्य करतात.

%%K t09-c1-15-2 p
\VUgras{मैलाचे दगड} \textsuperscript{(36)} आणि \VUgras{दिशादर्शक खांब} \textsuperscript{(37)} त्या वाटा दाखवतात, ज्यांवर अनेकदा एखादा तरुण \VUgras{डोंगरवासी} \textsuperscript{(38)} भेटतो — कमरेला \VUgras{झोळी} \textsuperscript{(40)} आणि हातात \VUgras{मारमॉट} \textsuperscript{(39)} घेतलेला —, किंवा एखादा \VUgras{बंजारा} \textsuperscript{(41)}, ज्याचा \VUgras{कोट} \textsuperscript{(42)} त्याच्या जुन्या फाटक्या \VUgras{टोपीशी} \textsuperscript{(43)} विचित्र विरोध करतो. तो साखळीने एक शिकवलेले \VUgras{अस्वल} \textsuperscript{(44)} फिरवतो.

%%K t09-tit-7 sub
\VUpk{10.2pt}{40}{डोंगरचढाई.}
\VUsaut{3.0mm}

%%K t09-c1-16-1 p
१६. --- जवळजवळ रोज मी \VUgras{वाटाड्याबरोबर} \textsuperscript{(48)} \VUgras{चढाईचा} सराव
करायला जातो, आणि पाठीला \VUgras{पिशवी} \textsuperscript{(47)} आणि हातात « \VUgras{अल्पेनस्टॉक}
» घेऊन, खऱ्या \VUgras{पर्यटकासारखा} \textsuperscript{(46)} मी \VUgras{खडकांवर}
\textsuperscript{(45)} चाल करतो तेव्हा मला फार अभिमान वाटतो.

%%K t09-c1-17-1 p
१७. --- डोंगराच्या माथ्यावरून सपाटीवरची माणसे मुंग्यांहून मोठी दिसत नाहीत;
\VUgras{मेंढ्यांचा कळप} \textsuperscript{(50)}, \VUgras{कुरणात} \textsuperscript{(49)} चरणारा,
लहान मुलांचे खेळणे वाटतो. एखादे \VUgras{पाइनचे झाड} \textsuperscript{(51)} किंवा एखादे
\VUgras{लाकडी घर} \textsuperscript{(52)} हिरवाईवरचा नुसता डाग वाटते.

%%K t09-c1-18-1 p
१८. --- या भटकंतीत आम्हाला \VUgras{पायवाटेवर} \textsuperscript{(53)} अनेकदा एखादा
\VUgras{चामोआची शिकार करणारा} \textsuperscript{(54)} भेटतो, नुकत्याच मारलेल्या जनावराला
खांद्यावर घेऊन जाणारा.

%%K t09-c1-19-1 p
१९. --- माझ्या वनस्पतिसंग्रहात मी \VUgras{नेच्याची} \textsuperscript{(55)} एक फार लक्षणीय
फांदी आणि \VUgras{हीथची} \textsuperscript{(57)} एक सुंदर गुलाबी डहाळी ठेवली आहे, जी मी माझ्या
शेवटच्या फेरफटक्यात एका लहानशा \VUgras{गुहेच्या} \textsuperscript{(56)} तोंडाशी तोडली.

%%K t09-tit-8 sub
\VUsaut{3.0mm}
\VUcentre{12.0pt}{0}{\textit{दुसरा प्रवेश.}}
\VUsaut{3.0mm}

%%K t09-tit-9 sub
\VUpk{11.4pt}{40}{रान. --- शिकार.}
\VUsaut{3.0mm}

%%K t09-c2-01-1 p
१. --- काल मी रानात एक अत्यंत सुखाचा दिवस घालवला. तिथे राहणाऱ्या प्राण्यांमध्ये आणि
\VUgras{कीटकांमध्ये} \textsuperscript{(5)} चालणाऱ्या उद्योगाने मला फार गुंतवून ठेवले.

%%K t09-c2-01-2 p
\VUgras{कोळी} \textsuperscript{(1)}, जो जाणाऱ्या माशीला पकडायला दोन फांद्यांमध्ये आपले
\VUgras{जाळे} \textsuperscript{(2)} विणतो, ती \VUgras{माशी} \textsuperscript{(3)}, आपले शिंपले
कष्टाने वाहणारी \VUgras{गोगलगाय} \textsuperscript{(4)}, आपल्या भक्ष्याची चाचपणी करणारा
शिंगवाला भुंगा, \VUgras{वारुळाजवळ} \textsuperscript{(6)} अतिशय उत्साही असलेली कष्टाळू मुंगी —
ती सगळी लहानगी जाता-येता विलक्षण उद्योगाने काम करत होती.

%%K t09-c2-02-1 p
२. --- एक \VUgras{साप} \textsuperscript{(7)}, जो पहिल्या नजरेत मला \VUgras{फुरसे} वाटला, पण जो
साधी \VUgras{धामण} याखेरीज दुसरे काही नव्हता, तो एका झाडाच्या खोडाखाली दडून बसला होता, आणि
मधूनमधून आपले बारीक डोके बाहेर काढत होता, जे नेहमी चावायला तयार असल्यासारखे वाटे. माझ्यासमोर
एक मोठी \VUgras{बिनशिंपल्याची गोगलगाय} \textsuperscript{(8)} जड पावलांनी सरकत होती, तिथे मुबलक
उगवणाऱ्या चवदार \VUgras{रानस्ट्रॉबेरींपैकी} \textsuperscript{(11)} एक फस्त केल्यावर.

%%K t09-c2-03-1 p
३. --- जमिनीवर \VUgras{रानफुलांचा} गालिचा पसरला होता; \VUgras{प्रिम्रोझ}
\textsuperscript{(9)}, \VUgras{सदाफुली} \textsuperscript{(10)}, आणि मला माहीत नसलेली आणखी
कितीतरी झाडे \VUgras{गवताच्या झुबक्यांमध्ये} \textsuperscript{(12)} फुलली होती.

%%K t09-c2-04-1 p
४. --- त्या प्रदेशात \VUgras{ट्रफल} जवळजवळ \VUgras{अळंबीएवढेच} \textsuperscript{(14)} सामान्य
आहे, आणि वनरक्षकाच्या मालकीचे एक \VUgras{डुकराचे पिल्लू} \textsuperscript{(13)} त्यातली काही
आपल्याला सापडतील का याची चटोरपणे चाचपणी करत होते.

%%K t09-tit-10 sub
\VUsaut{3.0mm}
\VUpk{10.2pt}{40}{चोरटी शिकार.}
\VUsaut{3.0mm}

%%K t09-c2-05-1 p
५. --- मला फार मजा वाटलेल्या एका प्रसंगाचा \VUgras{शेवट} मी पाहिला. आपल्या
\VUgras{बुटक्या पायांच्या शिकारी कुत्र्याच्या} \textsuperscript{(15)} हुंगण्याच्या मदतीने
\VUgras{वनरक्षकाने} \textsuperscript{(16)} नुकताच एका \VUgras{चोरट्या शिकाऱ्याला}
\textsuperscript{(22)} रंगेहात पकडले होते — नेमक्या त्याच क्षणी, जेव्हा तो मारलेली
\VUgras{शिकार} \textsuperscript{(17)} घेऊन जायची तयारी करत होता. पकडले जाण्यापेक्षा आपले
भक्ष्य सोडणे पसंत करून तो चोरटा शिकारी पळून गेला होता, आणि \VUgras{ससा}
\textsuperscript{(18)}, \VUgras{हरिणी} \textsuperscript{(19)}, \VUgras{भेकर}
\textsuperscript{(20)} आणि \VUgras{रानडुक्कर} \textsuperscript{(21)} गवतावर पडून होते,
वनरक्षकाला आनंद देत.

%%K t09-c2-06-1 p
६. --- रानात असलेल्या झाडांची विविधता थक्क करणारी आहे. \VUgras{बीचची झाडे}
\textsuperscript{(23)} आणि \VUgras{भूर्ज} \textsuperscript{(26)}, \VUgras{पाइन}, जे
\VUgras{राळ} \textsuperscript{(27)} देतात, \VUgras{ओक} \textsuperscript{(29)} आणि आणखीही
कितीतरी झाडे आपल्या फांद्या एकमेकांत मिसळतात.

%%K t09-c2-06-2 p
उंच झाडांच्या खोडांना चिकटणारा \VUgras{आयव्ही} \textsuperscript{(24)} \VUgras{रानाला}
\textsuperscript{(25)} चमकत्या हिरवाईने रंगवतो, जी \VUgras{शेवाळाच्या} \textsuperscript{(31)}
अधिक फिकट आणि निस्तेज हिरव्या छटेशी सुखद रीतीने जुळते — त्या शेवाळात \VUgras{सुतारपक्षी}
\textsuperscript{(30)} कीटक शोधतो.

%%K t09-tit-11 sub
\VUsaut{3.0mm}
\VUpk{10.2pt}{40}{रानातले निसर्गदृश्य.}
\VUsaut{3.0mm}

%%K t09-c2-07-1 p
७. --- जुन्या ओकनी मला भुरळ घातली. त्यांची जाड पिळवटलेली \VUgras{मुळे} \textsuperscript{(32)},
अर्धी जमिनीबाहेर आलेली, त्यांना बळाचे आणि जोमाचे भव्य रूप देतात, आणि मी स्वतःला विचारतो की
\VUgras{मालक} \textsuperscript{(35)} ती तोडून \VUgras{करवतीच्या} \textsuperscript{(34)} हवाली
करायला कसा तयार होतो — \VUgras{लाकूडतोड्याच्या} \textsuperscript{(33)} करवतीच्या. बिचारी
\VUgras{खोडे} \textsuperscript{(36)}, आपल्या \VUgras{सालीपासून} \textsuperscript{(37)} वंचित
झालेली किंवा \VUgras{कुऱ्हाडीने} \textsuperscript{(38)} चिरलेली — \VUgras{रोजंदार मजुराच्या}
\textsuperscript{(39)} —, मला दुःखी करतात.

%%K t09-c2-08-1 p
८. --- जवळच एका म्हाताऱ्या \VUgras{कोळसा भाजणाऱ्याने} \textsuperscript{(41)} आपल्या
\VUgras{फावड्याने कोळशाचा ढीग} \textsuperscript{(42)} तयार केला होता, आणि एका म्हातारीने
तोडलेल्या झाडांच्या डहाळ्यांच्या \VUgras{सुक्या लाकडाची} \VUgras{मोळी} \textsuperscript{(40)}
बांधून नेली.

%%K t09-tit-12 sub
\VUsaut{3.0mm}
\VUpk{10.2pt}{40}{शिकार.}
\VUsaut{3.0mm}

%%K t09-c2-09-1 p
९. --- काही क्षण विश्रांती घ्यायला मी \VUgras{वनरक्षकाच्या घरी} \textsuperscript{(46)} जायचे
ठरवले, पण \VUgras{तळ्याजवळ} \textsuperscript{(43)} पोहोचलो — ज्यावर एक सुंदर \VUgras{राजहंस}
\textsuperscript{(45)} पोहतो — तेव्हा माझ्या लक्षात आले की अलीकडच्या \VUgras{पुराने}
\textsuperscript{(44)} वाटेची खरीखुरी \VUgras{दलदल} करून टाकली आहे. मला वळसा घालावा लागला, आणि
रानाच्या कडेला असलेल्या \VUgras{देवदाराजवळून} \textsuperscript{(49)}, \VUgras{चिनारांजवळून}
\textsuperscript{(48)} आणि \VUgras{एल्मजवळून} \textsuperscript{(47)} जाताना मला
\VUgras{झुडपी रानातून} \textsuperscript{(54)} एक देखणे \VUgras{सांबर} \textsuperscript{(50)}
बाहेर पडताना दिसले, ज्याचा हट्टी \VUgras{शिकारी कुत्र्यांचा कळप} \textsuperscript{(51)} पाठलाग
करत होता — ज्यांना \VUgras{शिंगानेही} \textsuperscript{(53)} चेतवले होते,
\VUgras{घोड्यावरच्या शिकाऱ्यांच्या} \textsuperscript{(52)}. ते बिचारे जनावर पार थकून गेले
होते. माझ्यापासून काही पावलांवर ते मरून पडले.

%%K t09-c2-10-1 p
१०. --- \VUgras{शिकारीच्या धावपळीच्या} गोंगाटाने ओढला गेलेला वनरक्षकाचा मुलगा घराबाहेर आला आणि
आम्ही दोघे पुन्हा रानात गेलो. त्याने एका जुन्या ओकच्या फांदीवर एक \VUgras{मॅगपाय}
\textsuperscript{(56)} पाहिला होता. तिचे घरटे \VUgras{पालवीत} \textsuperscript{(58)} लपले
असावे असे त्याला वाटले, आणि त्याला ते घरटे काढायचे होते.

%%K t09-c2-10-2 p
\VUgras{खारीसारखा} \textsuperscript{(61)} चपळ, तो \VUgras{फांदीवरून} \textsuperscript{(55)}
फांदीवर चढला, पण त्याला काहीच सापडले नाही आणि एका म्हाताऱ्या \VUgras{कावळ्याला}
\textsuperscript{(60)} उडवण्यापलीकडे काही जमले नाही — \VUgras{डहाळीवर} \textsuperscript{(57)}
बसलेल्या.

\VUsaut{4.0mm}
\VUfilet{22mm}
